\par\noindent\textbf{(1)} Gauss定律给 \(E_r(r)=\lambda/(2\pi\varepsilon r)\)，Ampère环路给 \(H_\phi(r)=I/(2\pi r)\)，仅在 \(a<r<b\) 适用。两导体间电压 \(V=\int_a^bE_rdr=\lambda\ln(b/a)/(2\pi\varepsilon)\)，必须明确符号随内导体正电荷方向决定。
\par\noindent\textbf{(2)} Poynting矢量为 \(\boldsymbol S=\boldsymbol E\times\boldsymbol H=E_rH_\phi\hat z\)。横截面积分 \(P=\int_a^bE_rH_\phi\,2\pi r\,dr=\lambda I\ln(b/a)/(2\pi\varepsilon)=VI\)，说明无耗散电缆中的能量主要经介电间隙而不是导体内部传播。
\par\noindent\textbf{(3)} 电能 \(U_E'=\int_a^b\varepsilon E_r^2\,\pi r\,dr=\lambda^2\ln(b/a)/(4\pi\varepsilon)\)，磁能 \(U_B'=\mu I^2\ln(b/a)/(4\pi)\)。故 \(C'=2\pi\varepsilon/\ln(b/a)\)，\(L'=\mu\ln(b/a)/(2\pi)\)，相速度 \(v=1/\sqrt{L'C'}=1/\sqrt{\mu\varepsilon}\)。
\par\noindent\textbf{(4)} Abraham型电磁动量密度可写 \(\boldsymbol g=\boldsymbol S/c_m^2\)，其中 \(c_m=1/\sqrt{\mu\varepsilon}\) 在理想无色散介质近似下成立，故 \(\boldsymbol g=\mu\varepsilon E_rH_\phi\hat z\)。若需处理介质机械动量及Abraham–Minkowski差异，必须指定封闭系统的应力—能量张量；不能将此局部表达当作普适总动量。
